\subsection{Ветвление (Условные операторы)}

Условные операторы \texttt{if}, \texttt{elif} и \texttt{else} позволяют программе принимать решения в зависимости от условий.

\begin{codefile}[python]{python/examples/01_python_basics/05_branching.py}
\end{codefile}

Обратите внимание на отступы (4 пробела) — они обязательны в Python и обозначают блоки кода.

\begin{taskbox}[Светофор]
Напишите программу, которая принимает цвет светофора (\texttt{"red"}, \texttt{"yellow"}, \texttt{"green"}) и выводит действие (\texttt{"Стоп"}, \texttt{"Приготовиться"}, \texttt{"Ехать"}).
\end{taskbox}

\subsubsection*{Задачи для самостоятельного решения}

\begin{enumerate}
    \item Дано число \texttt{x = 10}. Если оно больше 0, выведите "Положительное".
    \item Дано \texttt{n = 7}. Если делится на 2 без остатка — "Четное", иначе — "Нечетное".
    \item Дано \texttt{age = 15}. Если \texttt{age >= 18}, выведите "Взрослый", иначе "Ребенок".
    \item Дано \texttt{password = "1234"}. Если пароль совпадает с "secret", пишем "Доступ разрешен", иначе "Отказ".
    \item Даны \texttt{a = 10}, \texttt{b = 20}. Выведите большее из них.
    \item Дано \texttt{x = 0}. Используя \texttt{elif}, выведите "Плюс", "Минус" или "Ноль".
    \item Дано \texttt{color = "red"}. Если "green" — "Ехать", "yellow" — "Ждать", "red" — "Стоять".
    \item \texttt{voltage = 3.6}. Если < 3.5 — "Посадка!", иначе "Полет нормальный".
    \item Делится ли число 15 на 3 и на 5 одновременно?
    \item Дано \texttt{name = "Alex"}. Если длина имени > 5 букв, напишите "Длинное", иначе "Короткое".
    \item Дана строка \texttt{s = ""}. Если она пустая (длина 0), напишите "Пусто".
    \item \texttt{weight = 80}. < 50 — "Легкий", 50-100 — "Средний", > 100 — "Тяжелый".
    \item \texttt{armed = True}, \texttt{ready = False}. Если оба \texttt{True} — "Взлет", иначе "Ждать".
    \item \texttt{login = "admin"}. Если не равен "admin", вывести "Ошибка".
    \item \texttt{a=5, b=2, c=8}. Найдите самое большое число из трех через \texttt{if}.
\end{enumerate}
